\section{Side-Channel Attacks dan Countermeasures}

\subsection{Pengantar Side-Channel Attacks}

Side-channel attacks mengeksploitasi informasi yang bocor dari implementasi fisik kriptografi, bukan kelemahan matematis algoritma itu sendiri \cite{kochert1996}. Informasi ini dapat berupa waktu eksekusi, konsumsi daya, emisi elektromagnetik, atau suara.

\textbf{Jenis Side-Channel Attacks}:
\begin{itemize}
  \item \textbf{Timing Attacks}: Analisis variasi waktu eksekusi
  \item \textbf{Power Analysis}: Analisis konsumsi daya
  \item \textbf{Electromagnetic Analysis}: Analisis emisi EM
  \item \textbf{Acoustic Cryptanalysis}: Analisis suara
  \item \textbf{Cache Attacks}: Eksploitasi cache behavior
\end{itemize}

\subsection{Timing Attacks}

Timing attacks pertama kali didemonstrasikan oleh Paul Kocher pada tahun 1996 terhadap RSA \cite{kochert1996}. Prinsipnya: waktu eksekusi operasi kriptografi dapat bergantung pada data rahasia (kunci).

\subsubsection{Timing Attack pada RSA}

Dalam RSA, operasi modular eksponensial $c = m^e \bmod n$ menggunakan algoritma square-and-multiply yang waktu eksekusinya bergantung pada bit-bit kunci.

\begin{lstlisting}[style=cstyle, caption={Vulnerable RSA Implementation}]
// Vulnerable implementation - timing depends on key bits
uint64_t mod_exp_vulnerable(uint64_t base, uint64_t exp, uint64_t mod) {
    uint64_t result = 1;
    
    for (int i = 63; i >= 0; i--) {
        result = (result * result) % mod;  // Always executed
        
        if ((exp >> i) & 1) {  // Timing leak!
            result = (result * base) % mod;  // Only for key bits = 1
        }
    }
    
    return result;
}

// Timing attack measurement
void timing_attack_rsa() {
    uint64_t ciphertexts[1000];
    uint64_t timings[1000];
    
    // Generate random ciphertexts
    for (int i = 0; i < 1000; i++) {
        ciphertexts[i] = random_uint64();
    }
    
    // Measure timing for each ciphertext
    for (int i = 0; i < 1000; i++) {
        clock_t start = clock();
        mod_exp_vulnerable(ciphertexts[i], unknown_key, MODULUS);
        clock_t end = clock();
        
        timings[i] = end - start;
    }
    
    // Analyze timing differences to recover key bits
    analyze_timings(ciphertexts, timings, 1000);
}
\end{lstlisting}

\subsubsection{Countermeasures: Constant-Time Implementation}

\begin{lstlisting}[style=cstyle, caption={Constant-Time RSA Implementation}]
uint64_t mod_exp_constant_time(uint64_t base, uint64_t exp, uint64_t mod) {
    uint64_t result = 1;
    uint64_t temp = base % mod;
    
    for (int i = 63; i >= 0; i--) {
        // Always compute both possibilities
        uint64_t result_if_bit_0 = (result * result) % mod;
        uint64_t result_if_bit_1 = (result_if_bit_0 * temp) % mod;
        
        // Select based on key bit (constant-time selection)
        uint64_t mask = -((exp >> i) & 1);  // All 1s if bit=1, all 0s if bit=0
        result = result_if_bit_0 ^ (mask & (result_if_bit_1 ^ result_if_bit_0));
        
        // Always update temp
        temp = (temp * temp) % mod;
    }
    
    return result;
}

// Alternative: Always perform multiplication
uint64_t mod_exp_always_multiply(uint64_t base, uint64_t exp, uint64_t mod) {
    uint64_t result = 1;
    
    for (int i = 63; i >= 0; i--) {
        result = (result * result) % mod;
        
        // Always multiply, but conditionally apply result
        uint64_t temp = (result * base) % mod;
        uint64_t mask = (exp >> i) & 1;
        result = result + mask * (temp - result);
    }
    
    return result;
}
\end{lstlisting}

\subsection{Power Analysis}

Power analysis attacks mengukur konsumsi daya selama eksekusi operasi kriptografi \cite{kochert1999}. Ada dua jenis utama:

\begin{itemize}
  \item \textbf{Simple Power Analysis (SPA)}: Analisis satu trace daya
  \item \textbf{Differential Power Analysis (DPA)}: Analisis statistik multiple traces
\end{itemize}

\subsubsection{SPA pada DES}

SPA dapat mengeksploitasi perbedaan daya antara operasi yang berbeda dalam DES.

\begin{lstlisting}[style=cstyle, caption={DES Implementation Vulnerable to SPA}]
void des_round_vulnerable(uint32_t *L, uint32_t *R, uint32_t *K) {
    uint32_t expanded_R = expand_e(*R);
    uint32_t substituted_R = substitution_s(expanded_R, *K);
    uint32_t permuted_R = permutation_p(substituted_R);
    
    // SPA leak: conditional operation based on key bit
    if (K[0] & 1) {  // Key-dependent branch
        *L ^= permuted_R;
    } else {
        *L ^= permuted_R >> 1;  // Different operation
    }
}

// Power trace measurement and analysis
void spa_attack_des() {
    PowerTrace traces[100];
    
    // Collect power traces during DES encryption
    for (int i = 0; i < 100; i++) {
        traces[i] = measure_power_trace(des_encrypt, plaintext[i], key);
    }
    
    // Analyze traces to identify key-dependent operations
    analyze_power_traces(traces, 100);
}
\end{lstlisting}

\subsubsection{DPA Implementation}

DPA menggunakan statistik untuk mengekstrak informasi kunci dari banyak power traces.

\begin{lstlisting}[style=cstyle, caption={DPA Attack Framework}]
typedef struct {
    double samples[1000];  // Power trace samples
    uint8_t plaintext[8];   // Corresponding plaintext
} PowerTrace;

// DPA attack on DES S-box
void dpa_attack_des_sbox(PowerTrace *traces, int num_traces) {
    int hypothesis_counts[256][2] = {0};  // [key_hypothesis][bit_value]
    double hypothesis_sums[256][2] = {0};
    
    for (int trace_idx = 0; trace_idx < num_traces; trace_idx++) {
        PowerTrace *trace = &traces[trace_idx];
        
        // Try all possible key values for S-box input
        for (int key_hypothesis = 0; key_hypothesis < 256; key_hypothesis++) {
            // Compute intermediate value
            uint8_t intermediate = trace->plaintext[0] ^ key_hypothesis;
            uint8_t sbox_output = sbox1[intermediate >> 4][intermediate & 0x0F];
            int target_bit = (sbox_output >> 3) & 1;  // Target bit 3
            
            // Add power trace to appropriate hypothesis group
            hypothesis_counts[key_hypothesis][target_bit]++;
            hypothesis_sums[key_hypothesis][target_bit] += trace->samples[100];  // Sample at relevant time
        }
    }
    
    // Find key hypothesis with maximum correlation
    double max_correlation = 0;
    int best_key = 0;
    
    for (int key = 0; key < 256; key++) {
        double mean0 = hypothesis_sums[key][0] / hypothesis_counts[key][0];
        double mean1 = hypothesis_sums[key][1] / hypothesis_counts[key][1];
        double correlation = fabs(mean1 - mean0);  // Simple correlation measure
        
        if (correlation > max_correlation) {
            max_correlation = correlation;
            best_key = key;
        }
    }
    
    printf("Best key hypothesis: 0x%02x (correlation: %.4f)\n", best_key, max_correlation);
}
\end{lstlisting}

\subsection{Cache Attacks}

Cache attacks mengeksploitasi perilaku cache CPU untuk mengekstrak informasi kriptografi \cite{osvik2006}.

\subsubsection{Prime+Probe Attack}

Prime+Probe attack mengisi cache dengan data target, kemudian mengukur waktu akses untuk menentukan mana yang diakses oleh victim.

\begin{lstlisting}[style=cstyle, caption={Prime+Probe Attack Framework}]
#include <time.h>

#define CACHE_SIZE 64  // Simplified cache size
#define CACHE_LINE_SIZE 64

// Probe array for cache attack
char probe_array[CACHE_SIZE * CACHE_LINE_SIZE];

void prime_cache() {
    // Fill entire cache with probe array
    for (int i = 0; i < CACHE_SIZE; i++) {
        // Force cache line to be loaded
        probe_array[i * CACHE_LINE_SIZE] = i;
    }
}

void probe_cache(uint64_t *timings) {
    for (int i = 0; i < CACHE_SIZE; i++) {
        clock_t start = clock();
        volatile char dummy = probe_array[i * CACHE_LINE_SIZE];
        clock_t end = clock();
        
        timings[i] = end - start;
    }
}

// AES cache attack example
void aes_cache_attack() {
    uint64_t timings[CACHE_SIZE];
    
    // Prime cache
    prime_cache();
    
    // Victim encrypts data (accesses S-box entries)
    victim_aes_encrypt(plaintext, key);
    
    // Probe cache to find which entries were accessed
    probe_cache(timings);
    
    // Analyze timings to identify accessed cache lines
    for (int i = 0; i < CACHE_SIZE; i++) {
        if (timings[i] < THRESHOLD) {
            printf("Cache line %d was likely accessed\n", i);
        }
    }
}
\end{lstlisting}

\subsubsection{Flush+Reload Attack}

Flush+Reload attack menggunakan shared memory (shared libraries) untuk cache attacks.

\begin{lstlisting}[style=cstyle, caption={Flush+Reload Attack}]
#include <sys/mman.h>
#include <fcntl.h>

void flush_reload_attack() {
    // Open shared library (e.g., OpenSSL crypto library)
    int fd = open("/usr/lib/x86_64-linux-gnu/libcrypto.so.1.1", O_RDONLY);
    void *crypto_lib = mmap(NULL, 4096, PROT_READ, MAP_SHARED, fd, 0);
    
    // Find AES T-table address in shared library
    uint8_t *aes_ttable = find_aes_ttable(crypto_lib);
    
    while (1) {
        // Flush T-table from cache
        for (int i = 0; i < 256; i++) {
            _mm_clflush(&aes_ttable[i * 64]);  // CLFLUSH instruction
        }
        
        // Wait for victim to perform AES operation
        usleep(1000);
        
        // Reload and measure access times
        for (int i = 0; i < 256; i++) {
            clock_t start = clock();
            volatile uint8_t dummy = aes_ttable[i * 64];
            clock_t end = clock();
            
            if (end - start < CACHE_HIT_THRESHOLD) {
                printf("T-table entry %d was accessed\n", i);
            }
        }
    }
}
\end{lstlisting}

\subsection{Countermeasures}

\subsubsection{Hardware Countermeasures}

\begin{itemize}
  \item \textbf{Random Delays}: Menambahkan noise timing
  \item \textbf{Power Filtering}: Menstabilkan suplai daya
  \item \textbf{Shielding}: Mengurangi emisi elektromagnetik
  \item \textbf{Secure Hardware}: TPM, HSM, secure enclaves
\end{itemize}

\subsubsection{Software Countermeasures}

\begin{lstlisting}[style=cstyle, caption={Side-Channel Resistant Implementation}]
// Blinding technique for RSA
uint64_t rsa_blinded_decrypt(uint64_t ciphertext, uint64_t d, uint64_t n) {
    // Generate random blinding factor
    uint64_t r = random_uint64() % n;
    uint64_t r_inv = mod_inverse(r, n);
    
    // Blind the ciphertext
    uint64_t blinded_ct = (ciphertext * mod_pow(r, 65537, n)) % n;
    
    // Decrypt blinded ciphertext
    uint64_t blinded_pt = mod_pow(blinded_ct, d, n);
    
    // Unblind the result
    uint64_t plaintext = (blinded_pt * r_inv) % n;
    
    return plaintext;
}

// Randomized projective coordinates for ECC
typedef struct {
    uint32_t x, y, z;  // Projective coordinates (X:Y:Z)
} ECPointProjective;

ECPointProjective ecc_blinded_multiply(ECPointProjective P, uint32_t scalar) {
    ECPointProjective result = {0, 1, 0};  // Point at infinity
    
    // Add random point R
    ECPointProjective R = generate_random_point();
    ECPointProjective PR = point_add(P, R);
    
    // Compute (scalar * P) + (scalar * R)
    result = scalar_multiply(PR, scalar);
    
    // Subtract (scalar * R)
    ECPointProjective sR = scalar_multiply(R, scalar);
    result = point_subtract(result, sR);
    
    return result;
}

// Constant-time memory operations
void constant_time_memcpy(void *dest, const void *src, size_t n) {
    volatile uint8_t *d = (volatile uint8_t*)dest;
    const volatile uint8_t *s = (const volatile uint8_t*)src;
    
    for (size_t i = 0; i < n; i++) {
        d[i] = s[i];  // Prevent compiler optimization
    }
}

// Constant-time comparison
int constant_time_compare(const void *a, const void *b, size_t n) {
    const volatile uint8_t *aa = (const volatile uint8_t*)a;
    const volatile uint8_t *bb = (const volatile uint8_t*)b;
    
    uint8_t result = 0;
    for (size_t i = 0; i < n; i++) {
        result |= aa[i] ^ bb[i];
    }
    
    return result;  // 0 if equal, non-zero if different
}
\end{lstlisting}

\subsection{Praktikum: Side-Channel Analysis}

\begin{lstlisting}[style=cstyle, caption={Simple Timing Analysis}]
#include <time.h>
#include <stdio.h>

// Function with timing dependency
int vulnerable_function(int secret, int input) {
    int result = input;
    
    // Timing leak: loop iterations depend on secret
    for (int i = 0; i < secret; i++) {
        result = (result * 1103515245 + 12345) & 0x7fffffff;
    }
    
    return result;
}

// Timing analysis
void timing_analysis() {
    int inputs[100];
    clock_t timings[100];
    
    // Generate inputs
    for (int i = 0; i < 100; i++) {
        inputs[i] = i;
    }
    
    // Measure timing for different secret values
    for (int secret = 1; secret <= 10; secret++) {
        clock_t total_time = 0;
        
        for (int i = 0; i < 100; i++) {
            clock_t start = clock();
            vulnerable_function(secret, inputs[i]);
            clock_t end = clock();
            
            total_time += (end - start);
        }
        
        printf("Secret %d: Average time = %ld ticks\n", secret, total_time / 100);
    }
}

int main() {
    timing_analysis();
    return 0;
}
\end{lstlisting}

\subsection{Evaluasi Keamanan Implementasi}

\textbf{Checklist Keamanan Side-Channel}:
\begin{itemize}
  \item \textbf{Timing}: Apakah waktu eksekusi konstan terhadap data rahasia?
  \item \textbf{Power}: Apakah konsumsi daya tidak tergantung pada operasi?
  \item \textbf{Cache}: Apakah akses memori tidak membocorkan informasi kunci?
  \item \textbf{Branching}: Apakah tidak ada conditional branches berdasarkan data rahasia?
  \item \textbf{Memory}: Apakah tidak ada data rahasia yang tersimpan dalam memori yang tidak perlu?
\end{itemize}

\begin{catatan}
Side-channel attacks menunjukkan bahwa keamanan kriptografi tidak hanya bergantung pada matematika, tetapi juga pada implementasi fisik. Countermeasures yang efektif memerlukan pendekatan holistik yang mencakup hardware, software, dan prosedur operasional.
\end{catatan}
